\section{KI-gest\"utzte Lerninhaltsgenerierung}
\label{sec:aigeneration}

Neben der RAG-basierten Chat-Funktionalit\"at bietet EduShelf zwei weitere KI-Features: die automatische Generierung von \textbf{Lernkarten (Flashcards)} und \textbf{Multiple-Choice-Quizzes}. Beide Features nutzen denselben Ollama-LLM-Stack, folgen aber einem eigenst\"andigen Workflow, hierbei wird keine Vektorsuche ben\"otigt, da der Nutzer direkt den Kontext  \"uber Dokumenten-IDs oder RAG-abgerufene Chunks liefert.

\subsection{Generierungspipeline (5 Schritte)}

Beide Services (\texttt{FlashcardService}, \texttt{QuizService}) folgen derselben f\"unfstufigen Pipeline:

\begin{enumerate}
    \item \textbf{Kontextbeschaffung:} Der Service akzeptiert drei Eingabequellen: einen direkten \texttt{Context}-String (aus der RAG-Pipeline), eine einzelne \texttt{DocumentId} oder eine Liste von \texttt{DocumentIds}. Im letzteren Fall werden die Textinhalte aller Dokumente via \texttt{GetDocumentContentAsync} abgerufen und mit einem Trennzeichen zusammengef\"ugt.
    \item \textbf{Kontext-Truncation:} Um das Kontextfenster des LLMs nicht zu sprengen, wird der Dokumenteninhalt auf maximal 12.000 Zeichen (Flashcards) bzw. 15.000 Zeichen (Quizze) begrenzt.
    \item \textbf{Prompt-Vorbereitung:} Der System-Prompt wird aus \texttt{appsettings.json} geladen. Die Variable \texttt{\{Count\}} wird durch die gew\"unschte Anzahl ersetzt.
    \item \textbf{LLM-Inferenz:} Das Modell wird mit dem vorbereiteten \texttt{ChatHistory}-Objekt aufgerufen.
    \item \textbf{JSON-Parsing und Persistenz:} Die Antwort wird geparst und in der Datenbank gespeichert.
\end{enumerate}

\subsection{Flashcard-Generierung (\texttt{FlashcardService})}

Der LLM wird angewiesen, eine JSON-Liste von Objekten mit den Feldern \texttt{Question} und \texttt{Answer} zur\"uckzugeben. Das folgende Snippet zeigt den Prompt-Aufbau und die anschlie\ss{}ende Verarbeitung:

\begin{lstlisting}[language={[Sharp]C}, caption={KI-Aufruf und JSON-Parsing in FlashcardService.cs}]
// Prompt zusammenbauen und KI aufrufen
var chatHistory = new ChatHistory();
chatHistory.AddSystemMessage(
    promptTemplate.Replace("{Count}", request.Count.ToString()));
chatHistory.AddUserMessage(
    $"Text to analyze:\n\n{documentContent}\n\n" +
    "IMPORTANT: Output strictly valid JSON. Start with [.");

var result = await chatCompletionService
    .GetChatMessageContentAsync(chatHistory);

// JSON bereinigen und deserialisieren
var cleanedJson = JsonHelper.ExtractJson(result.Content);
var flashcards = JsonSerializer.Deserialize<List<GeneratedFlashcardJson>>(
    cleanedJson,
    new JsonSerializerOptions { PropertyNameCaseInsensitive = true });
\end{lstlisting}

Nach dem Parsing werden die generierten Karten automatisch mit einem Tag versehen, der den Dokumenttitel tr\"agt, und als \texttt{Flashcard}-Entit\"aten in der Datenbank gespeichert.

\newpage
\subsection{Quiz-Generierung (\texttt{QuizService})}

Die Quiz-Generierung ist komplex, da das LLM ein verschachteltes JSON-Objekt mit \texttt{Title}, \texttt{Questions} und pro Frage mehrere \texttt{Answers} (jeweils mit \texttt{IsCorrect}-Flag) produzieren muss. Da LLMs trotz strikter Anweisung gelegentlich vom erwarteten Schema abweichen (z.B. alternative Feldnamen wie \texttt{options} statt \texttt{Answers}), implementiert der \texttt{QuizService} eine \textbf{zweistufige Parsing-Strategie}:

\begin{lstlisting}[language={[Sharp]C}, caption={Robustes JSON-Parsing in QuizService.cs}]
var jsonNode = JsonNode.Parse(cleanedJson);

// Stufe 1: Wrapper-Handling (falls LLM { "quiz": {...} } liefert)
if (jsonNode["quiz"] != null)
    jsonNode = jsonNode["quiz"];

// Stufe 2a: Standard-Deserialisierung versuchen
var result = jsonNode.Deserialize<GeneratedQuizJson>(options);

// Stufe 2b: Fallback -- Manuelles Parsing bei alternativen Schemas
if (result == null || !result.Questions.Any())
{
    // Felder wie "text"/"Text", "options"/"answers", "answer" (string)
    // werden manuell ausgelesen und normalisiert.
    result = ManuallyParseQuizJson(jsonNode);
}
\end{lstlisting}

Die manuelle Fallback-Stufe erkennt dabei sowohl Antwortoptionen als einfache Strings (\texttt{["Paris", "London", ...]}) als auch als Objekte (\texttt{[\{"text": ..., "isCorrect": ...\}]}). Dies macht die Quiz-Generierung robust gegen\"uber variierenden LLM-Antwortformaten.

Nach erfolgreichem Parsing wird das Quiz als \texttt{Quiz}-Entit\"at mit allen \texttt{Question}-- und \texttt{Answer}-Objekten in einer einzigen EF Core-Transaktion gespeichert.
